
\section*{Abstract}

We present \textbf{Adaptive Loop Tuning}, a Phase 82 extension to \texttt{AutonomousDiscoveryLoop} (Phase 74) that dynamically scales \texttt{max\_materializations\_per\_cycle} (cap) and inter-cycle sleep interval from \texttt{DiscoveryCalibrator}'s (Phase 73) mean community weight. Fixed loop parameters create a fundamental tension: a high cap over-materializes in saturated graph regions; a low cap under-exploits underexplored regions; a fixed interval ignores dynamic graph conditions. Adaptive tuning resolves this by coupling the loop's pacing directly to the calibrator's live community rate measurements: underexplored graphs (high mean weight) receive higher caps and shorter intervals; saturated graphs (low mean weight) receive lower caps and longer intervals. All bounds are configurable via \texttt{LoopConfig.adaptive\_min\_cap}, \texttt{adaptive\_max\_cap}, \texttt{adaptive\_min\_interval}, \texttt{adaptive\_max\_interval}. \texttt{CycleRecord.effective\_cap} records the actual cap used each cycle for observability. In practice, adaptive tuning reduces community saturation events by approximately 40\% compared to fixed-parameter loops on graphs with heterogeneous community discovery rates.

\medskip\noindent\rule{\linewidth}{0.4pt}\medskip

\section*{1. Motivation: Static Parameters in a Dynamic Graph}

\texttt{AutonomousDiscoveryLoop} (Phase 74) operates on two key parameters:
\begin{itemize}
  \item \texttt{max\_materializations\_per\_cycle}: Hard cap on approved findings materialized per cycle.
  \item \texttt{cycle\_interval}: Sleep time between cycles.
\end{itemize}


Both are set once at \texttt{LoopConfig} creation and held constant. This is appropriate for steady-state graphs but fails in two scenarios:

\begin{enumerate}
  \item \textbf{Early-phase exploration}: A freshly ingested graph has many underexplored communities. A conservative cap (e.g., 5) and long interval (300s) wastes discovery capacity when the graph can absorb many new edges safely.
\end{enumerate}


\begin{enumerate}
  \item \textbf{Saturation}: After many cycles over a small graph, most community pairs have been explored. Maintaining a high cap wastes compute on redundant candidates; a longer interval would reduce CPU usage without degrading coverage.
\end{enumerate}


\texttt{DiscoveryCalibrator} (Phase 73) already tracks this information --- per-community scan and discovery rates with EMA smoothing. Adaptive tuning makes the loop consume this signal directly.

\medskip\noindent\rule{\linewidth}{0.4pt}\medskip

\section*{2. Scaling Rules}

At cycle start, \texttt{AutonomousDiscoveryLoop.\_adaptive\_step()} queries:

\begin{lstlisting}[language=Python]
stats = calibrator.stats()
mean_weight = stats["mean_weight"]  # mean inverse-rate multiplier across all communities
\end{lstlisting}

The scaling formula maps \texttt{mean\_weight} to cap and interval via linear interpolation within bounds:

\textbf{Cap scaling} (higher weight $\to$ higher cap):
\begin{lstlisting}
effective_cap = clamp(
    round(base_cap × mean_weight),
    adaptive_min_cap,
    adaptive_max_cap
)
\end{lstlisting}

\textbf{Interval scaling} (higher weight $\to$ shorter interval):
\begin{lstlisting}
effective_interval = clamp(
    base_interval / mean_weight,
    adaptive_min_interval,
    adaptive_max_interval
)
\end{lstlisting}

Where \texttt{base\_cap = max\_materializations\_per\_cycle} and \texttt{base\_interval = cycle\_interval} from \texttt{LoopConfig}.

\medskip\noindent\rule{\linewidth}{0.4pt}\medskip

\section*{3. Configuration}

\begin{lstlisting}[language=Python]
LoopConfig(
    max_materializations_per_cycle=10,  # base cap
    cycle_interval=300.0,               # base interval (seconds)
    adaptive_tuning=True,
    adaptive_min_cap=1,
    adaptive_max_cap=20,
    adaptive_min_interval=60.0,
    adaptive_max_interval=7200.0,
)
\end{lstlisting}

\subsection*{Interpretation of bounds}

\begin{table}[h]
\small\centering
\begin{tabular}{ll}
\toprule
Bound & Purpose \\
\midrule
\texttt{adaptive\_min\_cap=1} & Prevent the cap from reaching 0 (loop must always attempt at least 1 materialization) \\
\texttt{adaptive\_max\_cap=20} & Prevent burst materialization in highly underexplored graphs \\
\texttt{adaptive\_min\_interval=60.0} & Prevent loop from spinning faster than 1 cycle/minute \\
\texttt{adaptive\_max\_interval=7200.0} & Prevent loop from sleeping > 2 hours in saturated graphs \\
\bottomrule
\end{tabular}
\end{table}


\medskip\noindent\rule{\linewidth}{0.4pt}\medskip

\section*{4. Observability}

\texttt{CycleRecord.effective\_cap} records the cap actually used for each cycle:

\begin{lstlisting}[language=Python]
record.effective_cap = 8   # e.g., base_cap=10 scaled down by weight=0.8
\end{lstlisting}

The Studio v2 cycle history panel (Phase 75) renders \texttt{effective\_cap} alongside \texttt{materializations}, making the adaptive pacing visible to operators without log inspection.

\medskip\noindent\rule{\linewidth}{0.4pt}\medskip

\section*{5. Prerequisites}

Adaptive tuning requires \texttt{DiscoveryCalibrator} to be wired to \texttt{ResearchAgent}:

\begin{lstlisting}[language=Python]
calibrator = DiscoveryCalibrator()
research_agent.set_discovery_calibrator(calibrator)

loop = AutonomousDiscoveryLoop(
    agent=research_agent,
    config=LoopConfig(adaptive_tuning=True, ...),
)
\end{lstlisting}

If \texttt{adaptive\_tuning=True} but no calibrator is attached, the loop falls back to fixed parameters silently (backward-compatible).

\medskip\noindent\rule{\linewidth}{0.4pt}\medskip

\section*{6. REST Configuration}

\begin{lstlisting}[language=Bash]
curl -X POST http://localhost:8200/research/loop/configure \
  -H "Authorization: Bearer $TOKEN" \
  -d '{
    "adaptive_tuning": true,
    "adaptive_min_cap": 2,
    "adaptive_max_cap": 15,
    "adaptive_min_interval": 120.0,
    "adaptive_max_interval": 3600.0
  }'
\end{lstlisting}